% Generated draft: requires mathematical and source review.
\documentclass[11pt]{article}
\usepackage[margin=25mm]{geometry}
\usepackage{fontspec}
\setmainfont{Latin Modern Roman}
\usepackage{amsmath,amssymb,mathtools,unicode-math}
\setmathfont{Latin Modern Math}
\usepackage{enumitem,needspace,xurl,hyperref}
\setlength{\parindent}{0pt}
\setlength{\parskip}{5pt}
\setlength{\emergencystretch}{3em}
\newcommand{\N}{\mathbb{N}}
\newcommand{\Z}{\mathbb{Z}}
\newcommand{\Q}{\mathbb{Q}}
\newcommand{\R}{\mathbb{R}}
\newcommand{\C}{\mathbb{C}}
\newcommand{\F}{\mathbb{F}}
\newcommand{\A}{\mathbb{A}}
\newcommand{\PP}{\mathbb{P}}
\newcommand{\E}{\mathbb{E}}
\newcommand{\eps}{\varepsilon}
\newcommand{\dd}{\,\mathrm{d}}
\begin{document}
\section*{UnsolvedMath ID: 30001676}
\textbf{Source problem code:} catalogue entry
\subsection{Definitions and mathematical statement}
A standard Borel space is a measurable space isomorphic, as a measurable space, to a Polish space equipped with its Borel sigma-algebra. Under the usual complex convention, a \(C^*\)-algebra is a complex Banach algebra \(A\) with a conjugate-linear involution \(a\mapsto a^*\) satisfying \((ab)^*=b^*a^*\) and \(\lVert a^*a\rVert=\lVert a\rVert^2\). It is separable if its norm topology has a countable dense subset.

The supplied question is: \emph{Are all reasonable parametrizations of the class of separable \(C^*\)-algebras by elements of standard Borel spaces equivalent?}

To display the quantifiers without supplying missing hypotheses, let \(\mathscr P\) denote the intended collection of parametrizations. The question has the schematic form
\[
\forall P,Q\in\mathscr P,\qquad
\bigl(\operatorname{Reasonable}(P)\land\operatorname{Reasonable}(Q)\bigr)
\Longrightarrow \operatorname{Equivalent}(P,Q).
\]
Here the collection \(\mathscr P\) and the predicates \(\operatorname{Reasonable}\) and \(\operatorname{Equivalent}\) are placeholders for definitions absent from the supplied record, rather than defined mathematical conditions.

In particular, the record does not specify the measurable structure required of a parametrization, whether coverage is understood up to \(*\)-isomorphism, or the maps and preservation properties required for equivalence. It also does not specify whether the algebras are required to be unital or whether units must be preserved. Thus the source question is retained, but a fully determined mathematical assertion cannot be extracted from this record alone.
\subsection{Short English statement}
The question asks whether every pair of reasonable standard Borel parametrizations of separable \(C^*\)-algebras is equivalent. The supplied record leaves both ``reasonable'' and ``equivalent'' undefined, so the intended universal assertion requires clarification before mathematical review.
\subsection{Sources}
\begin{itemize}
\item Local registry: \texttt{data/unsolvedmath/problems.json}, permanent UnsolvedMath ID \(30001676\), source problem code \texttt{catalogue entry}, title ``Parametrizations of Separable \(C^*\)-Algebras.''
\item Supplied report citation: \emph{Set Theory} (2011), \emph{Oberwolfach Reports}, pp.\ 85--140, DOI \texttt{10.4171/owr/2011/02}. Source URL and supplied literature reference labeled ``Oberwolfach problem statement'': \url{https://doi.org/10.4171/owr/2011/02}. The record attributes the question to the workshop-report collection; no individual proposer is supplied.
\item Supplied statement-verification metadata labels the statement \texttt{corrected\_verified} and reports that the official report was checked and the wording repaired to match its passage. This is a supplied provenance claim, not an independent verification performed for this draft.
\item Literature assessment dated exactly \texttt{2026-08-21}: the supplied curation labels the problem open and reports that no general resolution was found. Its listed checked source is the Oberwolfach report linked above. This assessment is supplied dated curation, not independently verified; the record expressly calls for independent verification of time-sensitive status and cited claims.
\end{itemize}
\end{document}
